\chapter{Reflection\\เปลี่ยนประสบการณ์เป็นความสามารถใหม่}

ประสบการณ์ไม่ได้สร้างการเรียนรู้โดยอัตโนมัติ การเรียนรู้เกิดเมื่อบุคคลหยุดตีความว่าเกิดอะไรขึ้น เหตุใดจึงเกิด และจะเปลี่ยนการตัดสินใจครั้งต่อไปอย่างไร Reflection จึงไม่ใช่ถ้อยคำปิดโครงการ แต่เป็นกลไกที่ทำให้ Portfolio และตัว Mentee เติบโต

\learningobjectives{
  \item ใช้ Reflection เพื่อแยกเหตุการณ์ การตีความ และการตัดสินใจ
  \item เปลี่ยนบทเรียนให้เป็นหลักการและการทดลองครั้งต่อไป
  \item สร้างวัฒนธรรมที่เรียนรู้จากทั้งความสำเร็จและความผิดพลาด
}

\section{วงจร What? So What? Now What?}
\textbf{What?} รวบรวมข้อเท็จจริงโดยไม่รีบตัดสิน: เกิดอะไรขึ้น ใครเกี่ยวข้อง และข้อมูลบอกอะไร \textbf{So What?} ตีความความหมาย: สมมติฐานใดถูกท้าทาย อะไรเป็นกลไกสำคัญ \textbf{Now What?} แปลงบทเรียนเป็นการเปลี่ยน Process การทดลอง หรือการตัดสินใจที่ระบุได้

\begin{examplebox}
\textbf{What?} ผู้ใช้ 7 จาก 10 คนหยุดใช้ต้นแบบหลังสัปดาห์แรก แม้ผลทดสอบทางเทคนิคดี\\
\textbf{So What?} ทีมออกแบบจาก workflow ในห้องปฏิบัติการและประเมินภาระงานจริงต่ำเกินไป\\
\textbf{Now What?} โครงการต่อไปจะสังเกต workflow ก่อนออกแบบ และใช้ retention สัปดาห์ที่สองเป็น decision gate ก่อนเพิ่ม feature
\end{examplebox}

\section{Reflection ที่ลึกกว่าการโทษ}
คำอธิบายว่า \enquote{ผู้ใช้ไม่พร้อม} หรือ \enquote{ทีมไม่ให้ความร่วมมือ} อาจมีส่วนจริง แต่ไม่สร้างอำนาจในการเปลี่ยนแปลง Mentor ควรถามต่อว่า ทีมออกแบบเงื่อนไขใดได้บ้าง สัญญาณใดถูกมองข้าม และ Mentee มีส่วนต่อสถานการณ์อย่างไร โดยรักษาความเมตตาต่อตนเองควบคู่กับความรับผิดชอบ

\section{บันทึกบทเรียนให้ใช้งานได้}
Reflection ที่ดีควรจบด้วยหลักการชั่วคราว เช่น \enquote{เมื่อผู้ใช้ต้องเปลี่ยน workflow ควรทดลองกับผู้ใช้จริงก่อนเลือกสถาปัตยกรรมระบบ} และระบุว่าจะทดสอบหลักการนั้นในโครงการใด ความรู้จึงเดินทางจากประสบการณ์หนึ่งไปยังบริบทใหม่ได้

\diagnostic{บทเรียนนี้เปลี่ยนสิ่งที่ Mentee จะทำในวันจันทร์หน้าอย่างไร}

\begin{mentornote}
จัด Reflection ก่อนที่ความทรงจำจะถูกทำให้เรียบเกินจริง หลัง milestone สำคัญให้กันเวลา 20 นาทีเพื่อบันทึกทั้งสิ่งที่เกิดตามคาด สิ่งที่ประหลาดใจ และคำถามที่ยังเปิดอยู่
\end{mentornote}

\begin{keytakeaway}
Reflection สมบูรณ์เมื่อบทเรียนเปลี่ยนวิธีมองหรือวิธีทำงานครั้งต่อไป ไม่ใช่เมื่อเขียนสรุปอดีตได้ไพเราะ
\end{keytakeaway}

\begin{reflection}
เลือกเหตุการณ์ที่ไม่เป็นไปตามคาด เขียน What, So What และ Now What อย่างละหนึ่งย่อหน้า จากนั้นระบุสมมติฐานหนึ่งข้อที่จะทดสอบ และสัญญาณที่จะบอกว่าบทเรียนนี้ใช้ได้จริง
\blankline\blankline
\end{reflection}

\chaptersummary{Reflection เชื่อมประสบการณ์กับการเติบโต ด้วยการแยกข้อเท็จจริง ตีความความหมาย และออกแบบการเปลี่ยนแปลงครั้งต่อไป เมื่อทำอย่างสม่ำเสมอ Reflection จะกลายเป็นความสามารถในการเรียนรู้ของทั้งบุคคลและทีม}
